\fsection{Descripción del proyecto}
El cliente es un centro logístico de comercio electrónico en el que los productos llegan a la zona de preparación
de pedidos amontonados y sin orden dentro de contenedores. Hoy en día el picking se hace a mano: un operario coge
cada pieza del contenedor y la deja en la bandeja de salida. Es un trabajo repetitivo y poco ergonómico, que se
convierte en un cuello de botella en los picos de demanda. Los sistemas robotizados convencionales que ha
valorado no le sirven, porque necesitan que las piezas lleguen ordenadas o en una posición conocida.

En la primera visita el cliente nos enseña los contenedores tal como llegan y nos plantea sus necesidades:
\begin{itemize}
	\item Automatizar la recogida de piezas directamente desde el contenedor, sin tener que ordenarlas ni
	      posicionarlas antes, porque los productos cambian con frecuencia y no quiere utillajes específicos
	      para cada uno.
	\item Que la solución pueda convivir con personas en el mismo espacio de trabajo, sin grandes vallados.
\end{itemize}

Para dar respuesta a estas necesidades, el presente proyecto desarrolla un sistema de selección automática de piezas
en entornos no estructurados, conocido como bin picking, combinando visión artificial 3D, inteligencia
artificial y robótica colaborativa. La solución se basa en la plataforma SIMATIC Robot Pick AI de Siemens,
una solución preentrenada que permite detectar objetos en un contenedor desordenado, calcular el punto óptimo
de agarre y guiar al robot para ejecutar la operación de forma autónoma, sin necesidad de posicionamiento
previo de las piezas. Al usarse un robot colaborativo, la celda puede compartir espacio con los operarios,
aunque se realiza la correspondiente evaluación de riesgos de la aplicación.

Cabe destacar que el proyecto parte del modelo de IA preentrenado proporcionado por Siemens, por lo que
el entrenamiento o reentrenamiento del modelo queda fuera del alcance de este trabajo. El foco se centra
en la integración, configuración y puesta en marcha del sistema completo.
